\subsection{Navigation}

The navigation stack of this robot, powered by nav2 and fields2cover, executes three main tasks:

\begin{itemize}
\item Mapping
\item Patrolling
\item Object Approach
\end{itemize}

\subsubsection{Mapping}

Before the robot can navigate properly through the environment and ensure proper coverage, this environment must first be known. Thats where mapping approaches come in handy. Several advanced and specialized mapping approaches exist already, but in this paper only three are considered. Frontier based mapping, \{2\}, \{3\}. These were chosen through a combination of novelty and frequent application in similar contexts. Once the environment is mapped \{sufficiently\} the robot is allowed to navigate the space and propagate its behavior further down the tree.

\begin{quote}
Add ros2 bag playback
\end{quote}

\subsubsection{Patrolling}

Two categories of coverage strategies are implemented:

\textbf{Systematic Navigation:} The environment is traversed using a structured pattern to ensure complete coverage. These approaches are often called coverage methods and paired with decomposition methods discussed in the theory section

\textbf{Reactive Navigation}: The robot dynamically traverses and selects exploration targets based on frontier regions or detected objects, prioritizing areas of high information gain.

\begin{quote}
iets meer over coverage planning.
\end{quote}

\textbf{Free space:}
The ROS framework allows for integration with a lot of modular software packages. throughout the years this has given rise to many universally applied packages, Nav2 is such a case.
Using Nav2 the Global Costmap can be used to extrapolate the free space for the robot to move in.